\section{Runtime and real-time operating points}
The runtime experiment measures the system that would actually be deployed: scene analysis, state update, detector execution at the selected resolution, confidence filtering, and tracker update. It uses a Tesla T4 and 200 frames. This is important for JRTIP because a quality improvement obtained by spending a uniformly larger detector budget is not equivalent to a controller that shapes compute over time.

For ByteTrack, v3 mean latency is 74.3238 ms, P95 latency 100.6847 ms, and throughput 13.4546 FPS. Detector time is 54.5688 ms, tracker time 7.9712 ms, and controller overhead 4.2562 ms. For OATrack, mean latency is 65.4748 ms, P95 85.8969 ms, and throughput 15.2731 FPS; detector, tracker, and controller components are 52.6399, 1.8355, and 4.2010 ms. The controller is therefore a small but nonzero part of each frame budget, roughly four milliseconds, while the detector remains dominant.

The baseline harness reports 139.93 ms and 7.15 FPS for ByteTrack and 128.85 ms and 7.76 FPS for OATrack at the declared high-resolution baseline. These values show the measured runtime reduction associated with the modern policy, but the harnesses do not include image-read work identically. We therefore do not state a perfectly component-matched speedup. The safe interpretation is that v3's selected lower-resolution actions reduce measured detector time while adding approximately 4.2 ms of explicit controller work.

Neither 13.4546 nor 15.2731 FPS should be described as 30-FPS real-time operation. They are measured T4 throughputs for a causal controller whose operating point varies over a sequence. The real-time contribution is the explicit quality--compute operating policy and its measured latency decomposition, not a claim that the system satisfies every application frame-rate target.

The runtime result also clarifies the role of the tracker. The difference between hosts is mainly tracker time: OATrack's measured tracker component is 1.8355 ms compared with ByteTrack's 7.9712 ms, while detector times are similar. This supports evaluating the same detector controller on multiple tracker hosts rather than attributing all throughput or quality behaviour to the controller alone. It also motivates reporting tracker-dependent quality and transfer effects separately.

The P95 values add a deployment perspective that a mean FPS number hides. ByteTrack's P95 is 100.6847 ms, compared with its 74.3238 ms mean; OATrack's P95 is 85.8969 ms, compared with a 65.4748 ms mean. The tail includes frames assigned to more expensive operating points and ordinary variability in the end-to-end path. A production system with a hard per-frame deadline would need a separate tail-latency policy or budget; the present freeze reports the measured distribution but does not claim deadline compliance.

The baseline comparison should also be read as an operating-point comparison rather than a detector-engineering benchmark. The frozen v3 actions include 1088 and 1280 pixel calls, so the dominant saving is the amount of detector computation selected over time. Scene analysis is intentionally simple and adds only about four milliseconds, but it is not free. Future cost-matched studies should compare schedules with equal mean detector work and identical image-read instrumentation if the objective is to isolate scheduling from average compute.
